\documentclass[a4paper,11pt]{jsarticle}
\usepackage[dvipdfmx,margin=24mm]{geometry}
\usepackage{amsmath,booktabs,array}
\usepackage[dvipdfmx,colorlinks=true,urlcolor=blue]{hyperref}
\title{エンドミル需要分布の動的予測\\\large 現行ガンマ法の検証と製造指示への適用}
\author{}\date{\today}
\begin{document}\maketitle
\begin{abstract}
本稿は、過去12か月の販売数量にガンマ分布を当てはめる現行の
製造指示方式について、需要変化への追随性を改善する方法を検討する。
ガンマ分布の形状・尺度、または平均・変動係数を時系列外挿する方法は
実行可能であるが、少数標本、移動窓の重複、ゼロ需要の扱いに課題がある。
そこで、現行法とパラメータ外挿法に加え、状態空間、回帰、カウント、
ハードルおよびCroston系モデルを候補とする。各モデルから製造指示間隔
$I$ とリードタイム $L$ の和に相当する累計需要分布を生成し、その分位点と
在庫ポジションから製造数量を決定する。モデルはADI--$\mathrm{CV}^2$
による需要分類と事業情報を用いて絞り込み、ローリング・オリジン評価と
在庫シミュレーションによって選択する。結論として、パラメータ外挿法を
移行しやすいベースラインとし、元の月次需要を直接扱うモデルを本命候補とする。
\end{abstract}

\section{はじめに}
\subsection{背景と目的}
型番別月次販売数量 $y_t$ は、型番切替、営業活動、用途消滅、季節性等で変化する。
本稿の目的は、時間変化を考慮した需要分布を予測し、定期見直し方式の
製造指示数量へ変換する一貫した方法を示すことである。具体的には、
次の三点を検討する。
\begin{enumerate}
\item 現行ガンマ法のパラメータ外挿は統計的・実務的に妥当か。
\item 需要特性に応じて、どの予測モデルを候補とすべきか。
\item 予測分布をどのように製造数量へ変換し、評価すべきか。
\end{enumerate}

製造判断が今後 $H$ 期間分を対象とするなら、一点の $y_{T+H}$ でなく累計
\[
D_{T,H}=\sum_{h=1}^{H}y_{T+h}
\]
の条件付き分布を予測する。本稿では、次回指示までの期間 $I$ と
リードタイム $L$ を合わせた保護期間 $H=I+L$ を用いる。

\subsection{対象データと前提}
基本データは型番別の月次販売数量とする。ただし販売数量は必ずしも
潜在需要ではない。欠品、受注制限、納期遅延がある月は販売が需要を
下回るため、受注残・失注・欠品記録による補正またはフラグ付けが必要である。
返品、取消、一時的な社内移動も通常需要と区別する。

$I$ と $L$ は月次データと同じ時間単位へ換算する。本稿では説明を明確に
するため整数期間かつ既知として扱う。リードタイム自体が変動する場合は、
需要だけでなく $L$ もシミュレーションする必要がある。価格、販促、
型番切替日等の外生変数は、予測時点で入手可能な値または計画値だけを用いる。

\section{現行ガンマ法の定式化と課題}
\subsection{分布とパラメータ推定}
ガンマ分布を形状 $k$、尺度 $\theta$ とすると
\[
f(x;k,\theta)=\frac{x^{k-1}e^{-x/\theta}}{\Gamma(k)\theta^k},
\quad \mathrm{E}[X]=k\theta,\quad \mathrm{Var}(X)=k\theta^2 .
\]
過去12か月の平均 $\bar y_t$ と、分母を $n$ とする分散 $s_t^2$ によるモーメント推定量は
\[
\hat k_t=\frac{\bar y_t^2}{s_t^2},\qquad
\hat\theta_t=\frac{s_t^2}{\bar y_t}.
\]
平均または分散が0なら定義できない。販売数量は0を含む整数なので、正の連続値だけを取る通常のガンマ分布は低需要品に適合しにくい。

\subsection{移動窓推定の限界}
12か月窓の情報中心は概ね5.5か月前であり、急変への反応が遅いという問題意識は妥当である。
しかし、推定値が厳密に約5.5か月前の真のパラメータを表すわけではない。
窓内で分布が変わると、推定分散にはランダムなばらつきだけでなくトレンドや季節性も混入する。

\section{現行法を拡張するパラメータ外挿}
\label{sec:parameter-extrapolation}
\subsection{\texorpdfstring{形状パラメータ $k$ と尺度パラメータ $\theta$}
{形状パラメータ k と尺度パラメータ theta}のモデル}

正値制約を保つモデル
\[
\log k_t=a_k+b_kt+S_{k,t}+u_{k,t},\qquad
\log\theta_t=a_\theta+b_\theta t+S_{\theta,t}+u_{\theta,t}
\]
を考える。各記号の意味は次のとおりである。
\begin{description}
\item[$t$] パラメータが表す月を示す時間変数である。12か月窓の終了月を
$r$ とするなら、その推定値は情報中心 $\tau_r=r-5.5$ に対応させる。
予測起点 $T$ における現在のパラメータは $t=T$、将来第 $h$ 期は
$t=T+h$ として求める。したがって、最後の窓推定値 $\tau_T=T-5.5$
から現在までの外挿距離は5.5期間、将来第 $h$ 期までは $h+5.5$ 期間となる。
\item[$a_k,a_\theta$] 時間変化と季節変動を除いた基準水準（切片）である。
これらは $k,\theta$ 自体の水準ではなく、その対数の水準を表す。
\item[$b_k,b_\theta$] 1か月当たりの長期的な増減傾向（傾き）である。
例えば $b_k>0$ なら、他の条件が同じとき形状パラメータ $k_t$ は
1か月ごとに約 $100b_k$\,\% 増加する。厳密な増加率は
$100\{\exp(b_k)-1\}$\,\% である。
\item[$S_{k,t},S_{\theta,t}$] 季節成分である。月次データなら通常は12か月周期とし、
各月固有の効果、正弦・余弦関数、または季節状態成分で表す。
切片と区別できるよう、1周期の季節効果の合計を0とする等の制約を置く。
\item[$u_{k,t},u_{\theta,t}$] モデルで説明できない短期変動と推定誤差である。
平均0の独立誤差とするほか、変化が翌月へ残る場合は自己回帰過程として扱う。
独立誤差なら将来予測の中心値では0と置くが、予測区間にはその不確実性を含める。
\end{description}

右辺をそれぞれ $\eta_{k,t},\eta_{\theta,t}$ と書けば、
\[
k_t=\exp(\eta_{k,t})>0,\qquad
\theta_t=\exp(\eta_{\theta,t})>0
\]
となる。指数関数は常に正なので、直線をそのまま外挿した場合に起こり得る
$k_t\leq0$ または $\theta_t\leq0$ を防げる。この意味で「正値制約を保つ」
モデルである。

\subsection{\texorpdfstring{平均 $\mu$ と変動係数 $c$}
{平均 mu と変動係数 c}のモデル}

ガンマ分布の平均と変動係数を
\[
\mu_t=\mathrm{E}[X_t]=k_t\theta_t,\qquad
c_t=\frac{\sqrt{\mathrm{Var}(X_t)}}{\mathrm{E}[X_t]}
    =\frac{1}{\sqrt{k_t}}
\]
と定義する。直近12か月のデータからは
\[
\hat\mu_t=\bar y_t,\qquad
\hat c_t=\frac{s_t}{\bar y_t}
\]
と推定できる。$\mu_t$ は1か月の平均需要量、$c_t$ は平均に対する
ばらつきの大きさを表す。例えば $c_t=0.5$ は標準偏差が平均の50\,\%、
$c_t=1$ は標準偏差と平均が同程度であることを意味する。

$\mu_t>0,\ c_t>0$ を保証するモデル関数案として
\[
\log\mu_t=a_\mu+b_\mu t+S_{\mu,t}+u_{\mu,t},\qquad
\log c_t=a_c+b_ct+S_{c,t}+u_{c,t}
\]
を用いる。記号の役割は前小節と同じであり、$a_\mu,a_c$ は対数尺度の
基準水準、$b_\mu,b_c$ は1か月当たりの傾向、$S_{\mu,t},S_{c,t}$ は
12か月周期等の季節成分、$u_{\mu,t},u_{c,t}$ は短期変動と推定誤差である。
例えば $b_\mu>0$ なら平均需要は増加傾向にあり、厳密な月次増加率は
$100\{\exp(b_\mu)-1\}$\,\% となる。$b_c>0$ なら、需要水準に対する
相対的なばらつきが時間とともに増えていることを示す。

\subsection{季節成分 $S_t$ と短期変動 $u_t$ のモデル候補}

前項の $S_t$ は規則的な季節変動、$u_t$ は傾向と季節性では説明できない
短期変動および推定誤差を表す。以下では添字 $\mu,c,k,\theta$ を省略して
一般形を示す。実際には対象とする各パラメータについて係数を個別に推定する。

\subsubsection{季節成分 $S_t$}

月別の形を柔軟に表す最も直接的な方法は、月別ダミーを用いる
\[
S_t=\sum_{j=2}^{12}\gamma_jD_{j,t}
\]
というモデルである。ここで $D_{j,t}$ は時点 $t$ が第 $j$ 月なら1、
それ以外なら0となる。12個すべての月効果を用いる場合は
\[
S_t=\gamma_{m(t)},\qquad \sum_{j=1}^{12}\gamma_j=0
\]
とし、切片との識別のために月効果の和を0に制約する。月ごとに任意の
季節効果を表現できて解釈しやすい反面、11個の自由パラメータが必要なので、
型番ごとの履歴が短い場合には過適合しやすい。

より少ないパラメータで滑らかな季節性を表すには、Fourier級数
\[
S_t=\sum_{h=1}^{H_F}\left\{
A_h\sin\left(\frac{2\pi ht}{12}\right)
+B_h\cos\left(\frac{2\pi ht}{12}\right)\right\}
\]
を用いる。$H_F=1$ なら年1回の周期を2係数で、$H_F=2$ なら半年周期を
含めて4係数で表せる。履歴が短い場合は月別ダミーより安定しやすいが、
特定月だけに生じる鋭いピークを表すには多数の調和項が必要となる。

季節性が年ごとに変化すると考える場合は、状態空間モデルとして
\[
S_t=S_{t-12}+\omega_t,\qquad
\omega_t\sim N(0,\sigma_S^2)
\]
とする確率的季節成分も考えられる。この場合も、例えば直近12か月について
$\sum_{j=0}^{11}S_{t-j}=0$ となる制約を置く。季節パターンの緩やかな変化を
追跡できる一方、推定する状態が増えるため、十分な履歴または複数型番間の
情報共有が必要である。

\subsubsection{短期変動 $u_t$}

最も単純な仮定は独立同分布の誤差
\[
u_t=\varepsilon_t,\qquad
\varepsilon_t\sim N(0,\sigma^2)
\]
である。このとき将来予測の対数尺度の中心では $u_{T+h}=0$ とするが、
予測区間には $\sigma^2$ を反映する。なお、例えば
$\log\mu_t=\eta_t+\varepsilon_t$ なら指数変換後の条件付き中央値は
$\exp(\eta_t)$ であるのに対し、条件付き平均は
\[
E(\mu_t\mid\eta_t)=\exp\left(\eta_t+\frac{\sigma^2}{2}\right)
\]
となる。したがって、製造数量の計算で平均を用いる場合は、この対数正規変換の
補正を無視しない。

一時的な変化が翌月以降にも残る場合は、AR(1)過程
\[
u_t=\rho u_{t-1}+\varepsilon_t,\qquad
|\rho|<1,\qquad \varepsilon_t\sim N(0,\sigma_\varepsilon^2)
\]
を用いる。この場合の $h$ か月先予測は
\[
\hat u_{T+h\mid T}=\rho^h\hat u_T
\]
であり、一時的な上振れまたは下振れが徐々に減衰する。さらに複雑な
自己相関が残る場合には
\[
u_t=\sum_{i=1}^{p_u}\phi_i u_{t-i}+\varepsilon_t
+\sum_{j=1}^{q_u}\vartheta_j\varepsilon_{t-j}
\]
というARMA$(p_u,q_u)$も候補となる。ただし型番ごとの系列が短い場合は、
高次数モデルを避け、独立誤差またはAR(1)を優先する。

需要水準が恒久的に変化すると考えるなら、ランダムウォーク
$u_t=u_{t-1}+\varepsilon_t$ も考えられる。しかしこれは短期誤差ではなく
時間変化する水準を表し、決定論的傾向 $bt$ と役割が重なる。その場合は
\[
\log\mu_t=\ell_t+S_t,\qquad
\ell_t=\ell_{t-1}+d_{t-1}+\varepsilon_{\ell,t},\qquad
d_t=d_{t-1}+\varepsilon_{d,t}
\]
という局所線形トレンドへ置き換え、水準 $\ell_t$ と傾き $d_t$ を状態として
逐次推定する方が明確である。

本データへの初期候補としては、パラメータ数を抑えた
\[
S_t=A_1\sin\left(\frac{2\pi t}{12}\right)
+B_1\cos\left(\frac{2\pi t}{12}\right),\qquad
u_t=\rho u_{t-1}+\varepsilon_t
\]
という一次Fourier季節性とAR(1)の組合せが考えられる。これを季節性なし、
独立誤差、月別ダミー等の単純な候補とローリング・オリジン評価で比較する。

ただし、12か月移動窓から計算した隣接するパラメータ推定値は11か月分の
観測を共有する。このため、推定された $u_t$ に強い自己相関があっても、
それが真の需要変化ではなく窓の重複によって機械的に生じた可能性がある。
残差の自己相関だけからAR次数を決定せず、窓の重複を再現したシミュレーション、
非重複窓による感度分析、および元の月次需要を直接扱う状態空間モデルとの
標本外比較によって確認する必要がある。

予測した $\mu_t,c_t$ は
\[
k_t=\frac{1}{c_t^2},\qquad
\theta_t=\mu_tc_t^2
\]
によりガンマ分布のパラメータへ戻せる。この変換では
$k_t\theta_t=\mu_t$ および
$\sqrt{k_t\theta_t^2}/(k_t\theta_t)=c_t$ が保たれる。
平均需要と相対的な不確実性を別々に説明できるため、$k,\theta$ を
直接予測する場合より、営業・製造担当者が結果を解釈しやすい。

データに明確な直線傾向がない場合は $b_\mu=b_c=0$ とした季節モデルを、
傾向が将来弱まると考える場合は、予測期間 $h$ に対して傾向の寄与を
$b(\phi+\phi^2+\cdots+\phi^h)$（$0<\phi<1$）とする減衰トレンドも候補にする。
型番切替や販促の時期が既知なら、その効果を表す説明変数を右辺に追加する。

\subsection{二つのモデル案に共通する注意点}

以上のモデルで現在および将来へ外挿することは可能だが、次の弱点がある。
\begin{enumerate}
\item 1回の推定が12点だけで、特に分散と $k$ が不安定である。
\item 隣接窓は11点を共有する。系列の滑らかさは窓の重複からも生じる。
\item $k,\theta$ を直接予測する場合、その誤差が平均 $k\theta$ と
分散 $k\theta^2$ で増幅し得る。
\item 線形外挿は型番切替等の変曲点に弱く、季節性には複数年の履歴が要る。
\item ゼロ需要と整数値を表現できない。
\end{enumerate}
両案をベースラインとしてバックテストし、解釈のしやすさに加えて、
$H=I+L$ 期間累計需要の予測精度と在庫成果により採用案を決める。

\section{候補となる需要予測モデル}
\label{sec:recommended-models}

\subsection{第\ref{sec:parameter-extrapolation}節との関係}

第\ref{sec:parameter-extrapolation}節と本節は、同じ目的に対する
二つの異なる予測経路を示している。関係を整理すると次のようになる。
\begin{center}
\begin{tabular}{>{\raggedright\arraybackslash}p{28mm}
                >{\raggedright\arraybackslash}p{50mm}
                >{\raggedright\arraybackslash}p{50mm}}
\toprule
 & 第\ref{sec:parameter-extrapolation}節 & 第\ref{sec:recommended-models}節\\
\midrule
位置付け & 現行ガンマ分布法の拡張 & 現行法を含めて比較する代替モデル群\\
予測する対象 & 移動窓から推定した $k,\theta$ または $\mu,c$
& 元の月次販売数量 $y_t$、または需要発生と正の需要量\\
予測の流れ & 販売実績 $\rightarrow$ パラメータ推定
$\rightarrow$ パラメータ外挿 $\rightarrow$ 需要分布
& 販売実績 $\rightarrow$ 将来需要分布\\
主な役割 & 問題意識に沿ったベースライン & 実運用で採用する本命候補\\
\bottomrule
\end{tabular}
\end{center}

したがって、第\ref{sec:parameter-extrapolation}節は本節と独立した
別目的の分析ではなく、本節で比較する候補の一つである。
パラメータ外挿法は現行業務から移行しやすい一方、元データから移動窓の
推定値を作ってから予測するため、推定誤差と情報の遅れを受ける。
これに対し、ETS、カウントモデル等は $y_t$ を直接予測するため、
より速く需要変化を反映できる可能性がある。優劣は理論だけで決めず、
同じ予測起点、同じ保護期間 $H=I+L$、同じ評価指標で比較する。

なお、両節は組み合わせることもできる。本節のハードルモデルでは、
需要発生確率 $p_t$ は別に予測し、需要が発生した場合の正の需要量に対して
第\ref{sec:parameter-extrapolation}節の $k,\theta$ モデルまたは
$\mu,c$ モデルを適用できる。この場合、第\ref{sec:parameter-extrapolation}節
はハードルモデルの「正の需要量」部分を構成する。

\subsection{各モデルの定式化}

\subsubsection{単純法}

単純法は比較の基準であり、モデル選択から外してはならない。予測起点を
$T$、予測先を $h=1,\ldots,H$ とすると、主な定義は次のとおりである。
\[
\begin{array}{ll}
\text{ナイーブ法:} & \hat y_{T+h|T}=y_T,\\[2mm]
\text{季節ナイーブ法:} & \hat y_{T+h|T}=y_{T+h-12},\\[2mm]
\text{$m$か月移動平均:} &
\displaystyle \hat y_{T+h|T}=\frac{1}{m}\sum_{j=0}^{m-1}y_{T-j}.
\end{array}
\]
季節ナイーブ法は前年同月の実績をそのまま用いるため、少なくとも2年以上の
履歴を必要とする。移動平均の窓幅 $m$ は3、6、12か月等を候補とし、
ローリング評価で選ぶ。点予測だけでなく分布が必要な場合は、過去の予測誤差を
月別または需要クラス別に再標本化し、予測値に加える残差ブートストラップを使う。
各月の模擬需要を合計すれば $H$ 期間累計需要の分布を作れる。ただし負の模擬値は
0に切り上げる。この方法は簡単で説明しやすい一方、傾向の変化や型番切替を
明示的に表せない。

\subsubsection{ETS・状態空間モデル}

ETSは誤差（Error）、傾向（Trend）、季節性（Seasonal）を平滑化しながら
逐次更新する。状態空間表現により尤度、予測区間、将来経路を一貫して
計算できる\cite{hyndman2002ets}。例として、減衰傾向を持つ加法モデルは
\begin{align}
\ell_t&=\alpha(y_t-s_{t-12})+(1-\alpha)(\ell_{t-1}+\phi b_{t-1}),\\
b_t&=\beta(\ell_t-\ell_{t-1})+(1-\beta)\phi b_{t-1},\\
s_t&=\gamma(y_t-\ell_t)+(1-\gamma)s_{t-12}
\end{align}
で表せる。$\ell_t$ は現在の需要水準、$b_t$ は傾向、$s_t$ は月別季節効果、
$0\leq\alpha,\beta,\gamma\leq1$ は更新の速さ、$0<\phi\leq1$ は傾向の
減衰率である。$h$ か月先の中心予測は
\[
\hat y_{T+h|T}=\ell_T+(\phi+\phi^2+\cdots+\phi^h)b_T+s_{T+h-12q}
\]
となる。ここで $q$ は季節添字を直近の同じ月へ戻す整数である。

水準・傾向・季節性の有無、加法・乗法、減衰の有無を候補とし、尤度と
標本外評価で選ぶ。状態空間表現を使えば、観測誤差と状態推定の不確実性を
伝播して予測区間を得られる。需要が大きく、ほぼ毎月発生する品目に向く。
加法モデルは負の予測を出し得るため、0に切り上げるか、
$\log(1+y_t)$ をモデル化する。ゼロが多い系列ではETS単独より
ハードル法またはCroston系を優先する。

\subsubsection{回帰付き時系列モデル}

販売数量と、需要変化の原因として利用できる変数を結び付ける。基本形は
\[
g\{\mathrm{E}(y_t\mid\boldsymbol{x}_t)\}
=\beta_0+\beta_1t+\sum_{m=2}^{12}\delta_mD_{m,t}
+\boldsymbol{\gamma}^{\mathsf T}\boldsymbol{x}_t+n_t
\]
である。$g$ はリンク関数、$D_{m,t}$ は月ダミー、
$\boldsymbol{x}_t$ は価格、販促、営業接触、顧客数、切替後フラグ等、
$n_t$ は必要に応じてARIMA等で表す自己相関誤差である。
連続量には恒等または対数リンク、カウントには対数リンクを用いる。

係数は最小二乗法または最尤法で推定する。$H$ 期間先まで予測するには
$\boldsymbol{x}_{T+1},\ldots,\boldsymbol{x}_{T+H}$ が必要なので、
計画として確定した価格・販促等だけを直接利用する。未確定の変数は
基準・上振れ・下振れのシナリオを作り、各シナリオの需要分布を示す。
実績集計後にしか分からない変数を使うと情報漏洩になる。

型番切替では旧品と新品を個別に予測すると合計需要を二重に見積もる恐れがある。
そこで、まずシリーズ全体の合計需要 $y_t^{\mathrm{old}}+y_t^{\mathrm{new}}$
を予測し、次に新品配分率
\[
r_t=\frac{y_t^{\mathrm{new}}}
{y_t^{\mathrm{old}}+y_t^{\mathrm{new}}},\qquad
\operatorname{logit}(r_t)=a_r+b_rt
\]
を予測して二つに配分する方法がある。切替日等の原因情報がある場合に強いが、
係数を増やし過ぎないよう正則化または品目間の階層化が必要である。

\subsubsection{カウントモデル}

販売数量の整数性と時系列依存を直接扱う\cite{davis2003count}。
Poissonモデルは
\[
y_t\sim\mathrm{Poisson}(\lambda_t),\qquad
\log\lambda_t=\beta_0+\beta_1t+S_t+
\boldsymbol{\gamma}^{\mathsf T}\boldsymbol{x}_t+
\rho\log(1+y_{t-1})
\]
とする。$\lambda_t$ は月次需要の条件付き平均であり、対数リンクにより
$\lambda_t>0$ が保証される。係数は最尤法で推定する。

Poisson分布は平均と分散がともに $\lambda_t$ という制約を持つ。
実データの分散が平均を大きく上回る場合は負の二項分布を使い、
\[
\mathrm{E}(y_t)=\lambda_t,\qquad
\mathrm{Var}(y_t)=\lambda_t+\frac{\lambda_t^2}{r}
\]
とする。$r>0$ は過分散の大きさを決めるパラメータで、$r$ が小さいほど
ばらつきが大きい。残差の自己相関が残る場合はラグ項または動的状態を追加する。

将来の説明変数と直前需要から $\lambda_{T+h}$ を順次計算し、
各月のPoissonまたは負の二項乱数を多数生成する。ラグ需要が未知になる
$h\geq2$ では、直前に生成した値を次月へ渡す。このシミュレーションから
各月と $H$ 期間累計の平均・分位点を得る。毎月需要が発生する低～中数量品には
適するが、構造的ゼロが多ければハードルモデルを用いる。

\subsubsection{ハードルモデル}

ハードルモデル\cite{mullahy1986}は「その月に需要が発生するか」と
「発生したとき何本売れるか」を二段階でモデル化する。発生部分は
\[
z_t=\mathbf{1}(y_t>0),\qquad
\operatorname{logit}(p_t)=a_p+b_pt+S_{p,t}
+\boldsymbol{\gamma}_p^{\mathsf T}\boldsymbol{x}_t,\qquad
z_t\sim\mathrm{Bernoulli}(p_t)
\]
とする。正の需要量部分は
\[
y_t\mid y_t>0\sim\mathrm{Gamma}(k_t,\theta_t)
\]
とし、$k_t,\theta_t$ には第\ref{sec:parameter-extrapolation}節の
モデルを用いる。平均・変動係数モデルを使うなら、正の需要量の平均を
$\mu_t^+$、変動係数を $c_t^+$ と表し、ゼロを除いた観測から推定する。

正の需要量の平均と分散を $\mu_t^+=k_t\theta_t$、
$(\sigma_t^+)^2=k_t\theta_t^2$ とすると、ゼロを含む需要全体では
\[
\mathrm{E}(y_t)=p_t\mu_t^+,\qquad
\mathrm{Var}(y_t)=p_t(\sigma_t^+)^2+
p_t(1-p_t)(\mu_t^+)^2
\]
となる。発生確率の不確実性も在庫計算へ反映できる点が利点である。

$H$ 期間予測では各月についてBernoulli乱数を生成し、1ならガンマ乱数、
0なら需要0とする。これを多数回繰り返して累計需要分布を作る。
販売数量が小さい整数で、正の数量にもガンマ分布が不自然なら、
正の部分をゼロ切断Poissonまたはゼロ切断負の二項分布へ置き換える。
二つのモデルを別々に推定するため、正の観測が少ない品目では
品目群で係数を共有するか、単純なTSBを優先する。

\subsubsection{Croston系（Croston、SBA、TSB）}

正の需要が発生した時点を $t_j$、その数量を $q_j=y_{t_j}$、
前回からの間隔を $a_j=t_j-t_{j-1}$ とする。Croston法は需要発生時だけ
\begin{align}
\hat q_j&=\alpha q_j+(1-\alpha)\hat q_{j-1},\\
\hat a_j&=\beta a_j+(1-\beta)\hat a_{j-1}
\end{align}
と更新し、1期間当たり需要を
\[
\hat y_t^{\mathrm{CRO}}=\frac{\hat q_j}{\hat a_j}
\]
と予測する。需要がない月には予測値を更新しない。
$\alpha,\beta$ は通常0から1の範囲を探索し、ローリング評価で選ぶ。
SBAは、需要間隔の平滑化係数を $\beta$ とした場合、
\[
\hat y_t^{\mathrm{SBA}}=
\left(1-\frac{\beta}{2}\right)\frac{\hat q_j}{\hat a_j}
\]
としてCroston法の正の偏りを近似的に補正する。

TSBは間隔でなく需要発生確率を平滑化する。各月で
\[
\hat p_t=\beta z_t+(1-\beta)\hat p_{t-1}
\]
と更新し、需要が発生した月だけ
\[
\hat q_t=\alpha y_t+(1-\alpha)\hat q_{t-1}
\]
と更新する。予測は $\hat y_t^{\mathrm{TSB}}=\hat p_t\hat q_t$ である。
ゼロの月にも $\hat p_t=(1-\beta)\hat p_{t-1}$ と下がるので、
終売や用途消滅にCroston法より速く反応する。

初期値は最初の数件の正の需要量・間隔、または学習期間全体の平均から置く。
初期値と平滑化係数の選択は結果に影響するため、全評価起点で同じ規則を使う。
Croston系は軽量で断続需要に適する一方、標準形は確率分布を直接与えず、
季節性や説明変数も扱いにくい。分布予測が必要なら、需要発生と数量の
経験分布を用いるシミュレーション、またはハードルモデルを併用する。

\subsection{予測分布から製造指示数量への変換}
\label{sec:production-quantity}

次回の製造指示までの期間を $I$、製造リードタイムを $L$ とし、
両者を需要データと同じ単位（例えば月）で表す。定期的に製造指示を出す
運用では、今回の指示で保護すべき期間は
\[
H=I+L
\]
である。この期間の累計需要を
\[
D_{T,H}=\sum_{h=1}^{H}y_{T+h}
\]
と定義し、どのモデルでも最終的に $D_{T,H}$ の予測分布を求めれば、
同じ製造指示の考え方を適用できる。

\paragraph{現行ガンマ法が成立する条件}
各期間の需要 $y_{T+h}$ が互いに独立で、すべて同じ
$\mathrm{Gamma}(k,\theta)$ に従うなら、ガンマ分布の加法性により
\[
D_{T,H}\sim\mathrm{Gamma}(Hk,\theta)
\]
となる。したがって、現在使用している $f((I+L)k,\theta)$ は、
この仮定の下では正しい。より一般に、形状 $k_h$ が期間ごとに異なっても
尺度が共通の $\theta$ なら
\[
D_{T,H}\sim\mathrm{Gamma}\left(\sum_{h=1}^{H}k_h,\theta\right)
\]
となる。しかし、第\ref{sec:recommended-models}節で扱うように
季節性や傾向によって $\theta_h$ も変わる場合、または期間間に相関がある
場合、累計需要は一般にはガンマ分布にならない。その場合に
単純に $Hk$ と置くと、需要の変化と相関を製造数量へ正しく反映できない。

\paragraph{各予測モデルでの計算}
各方法から保護期間需要を求める方法は次のとおりである。
\begin{center}
\begin{tabular}{>{\raggedright\arraybackslash}p{30mm}
                >{\raggedright\arraybackslash}p{102mm}}
\toprule
方法 & $D_{T,H}$ の予測分布の作り方\\
\midrule
単純法 &
各将来月の点予測に過去の予測誤差を再標本化して加え、$H$ 期間分を
合計する。季節別の誤差と期間間の相関を保つため、可能なら連続した
$H$ か月単位で残差を再標本化する。\\
ETS・状態空間 &
予測起点の状態から、観測誤差と状態誤差を含む将来経路を $H$ 期間生成し、
各経路を合計する。状態を共有するために生じる月間相関も反映できる。\\
回帰付き時系列 &
将来の説明変数または各シナリオを与え、回帰誤差と時系列誤差を生成して
$H$ 期間を合計する。説明変数が未確定ならシナリオ別の製造数量を示す。\\
カウントモデル &
各将来月のPoissonまたは負の二項分布から順番に乱数を生成し、合計する。
独立Poissonであれば累計もPoissonだが、動的な平均やラグがあれば
経路シミュレーションを用いる。\\
ハードルモデル &
各月について需要発生をBernoulli分布で生成し、発生時だけ正の需要量を
生成して合計する。ゼロ需要と需要量変化をともに反映できる。\\
Croston・SBA・TSB &
標準形は点予測しか与えないため、$\hat y H$ だけでは安全在庫を計算できない。
推定した発生確率・間隔と正の需要量の経験分布から需要経路を生成するか、
ハードルモデルで分布を補う。\\
\bottomrule
\end{tabular}
\end{center}

実装は共通化できる。各モデルから $M$ 本の将来需要経路
$y_{T+1}^{(m)},\ldots,y_{T+H}^{(m)}$ を生成し、
\[
D_{T,H}^{(m)}=\sum_{h=1}^{H}y_{T+h}^{(m)}
\qquad(m=1,\ldots,M)
\]
を計算する。この標本の平均を期待需要、分位点を安全在庫を含む
基準在庫量として用いる。例えば欠品しない確率の目標を
$\alpha=0.95$ とするなら、目標在庫ポジションは
\[
S_T=Q_\alpha(D_{T,H}\mid\mathcal{F}_T)
\]
である。$Q_\alpha$ は予測分布の $\alpha$ 分位点、
$\mathcal{F}_T$ は予測時点で利用可能な情報を表す。

現在庫を $OH_T$、入荷予定を含む発注・製造残を $OO_T$、
未充足需要を $BO_T$ とすると、在庫ポジションは
\[
IP_T=OH_T+OO_T-BO_T
\]
であり、今回の正味製造指示数量は
\[
q_T=\max(0,S_T-IP_T)
\]
となる。製造ロットを $B$ 本単位とするなら
\[
q_T^{\mathrm{lot}}
=B\left\lceil\frac{\max(0,S_T-IP_T)}{B}\right\rceil
\]
と切り上げる。既発行の製造指示を $OO_T$ に含めないと二重発注になるため、
在庫だけでなく在庫ポジションを用いることが重要である。

\paragraph{ガンマ分布の形式を維持する場合}
既存システムがガンマ分布しか受け付けない場合は、シミュレーションで得た
$D_{T,H}$ の平均 $m_H$ と分散 $v_H$ に対して
\[
k_H=\frac{m_H^2}{v_H},\qquad
\theta_H=\frac{v_H}{m_H}
\]
とモーメントを合わせ、$\mathrm{Gamma}(k_H,\theta_H)$ で近似できる。
ただしこれは一般には厳密な分布ではない。特にゼロ需要確率が高い場合や
分布が複数の山を持つ場合は分位点がずれるため、シミュレーション標本から
直接 $Q_\alpha$ を求める方法を優先する。近似を採用する場合は、
現行の $\mathrm{Gamma}(Hk,\theta)$ とともに、実績に対する分位点の
被覆率と在庫成果をバックテストする。

\section{需要のクラス分け}

\subsection{分類の計算}

評価期間を直近 $n$ か月、そのうち正の需要があった月数を $N$、
正の需要量を $q_1,\ldots,q_N$ とする。正の需要月同士の平均間隔をADI、
正の需要量の変動係数の二乗を $\mathrm{CV}^2$ とし、実務上は
\[
\mathrm{ADI}=\frac{n}{N},\qquad
\mathrm{CV}^2=\left(\frac{s_q}{\bar q}\right)^2
\]
で計算する。$\bar q$ と $s_q$ は正の需要量だけから求めた平均と標準偏差で
ある。ゼロ需要月を $\mathrm{CV}^2$ の計算に含めない点に注意する。
$N=0$ なら両指標を計算できず、$N=1$ なら $\mathrm{CV}^2$ を安定して
計算できないため、「実績なし」または「極少実績」として別管理する。

Syntetosら\cite{syntetos2005category}の目安は次のとおりである。
\begin{center}\begin{tabular}{lll}\toprule
分類 & 条件 & 第一候補\\\midrule
Smooth & ADI $\le1.32,\ \mathrm{CV}^2\le0.49$ & ETS、季節モデル\\
Erratic & ADI $\le1.32,\ \mathrm{CV}^2>0.49$ & 負の二項、ETS\\
Intermittent & ADI $>1.32,\ \mathrm{CV}^2\le0.49$ & SBA、ハードル\\
Lumpy & ADI $>1.32,\ \mathrm{CV}^2>0.49$ & TSB、ハードル・階層\\\bottomrule
\end{tabular}\end{center}
境界値は自然法則でなく特定手法の誤差比較から得た目安なので、自社データで
再検証する。分類期間は原則として直近24--36か月等に統一する。
全履歴を使うと現在は終売期の品目に過去の成長期が混ざり、短過ぎる窓では
偶然のゼロや大口注文でクラスが頻繁に変わる。ローリング評価では、
各予測起点までのデータだけでADIと $\mathrm{CV}^2$ を再計算する。

\subsection{分類の役割}

分類の目的は、四つのラベルから機械的に一つのモデルを決めることではない。
主に次の五つへ活用する。
\begin{enumerate}
\item 各品目について比較すべきモデル候補を絞り、計算量と保守負担を減らす。
\item 同じ需要特性を持つ品目をまとめ、モデルや平滑化係数を安定して推定する。
\item 予測精度と在庫成果をクラス別に集計し、全社平均に隠れる問題を発見する。
\item 重点的に人が確認すべき品目と、自動処理できる品目を分ける。
\item 需要構造の変化を検知し、モデル再選択や終売確認のきっかけにする。
\end{enumerate}

\subsection{モデル候補の絞り込み}

クラスごとの第一候補と実務上の確認事項は次のように設定する。
\begin{center}
\begin{tabular}{>{\raggedright\arraybackslash}p{22mm}
                >{\raggedright\arraybackslash}p{43mm}
                >{\raggedright\arraybackslash}p{67mm}}
\toprule
クラス & 比較する候補 & 運用上の確認\\
\midrule
Smooth &
季節ナイーブ、ETS、回帰付き時系列、現行ガンマ法 &
多くは自動予測しやすい。季節性、緩やかな傾向、継続的なバイアスを確認する。\\
Erratic &
移動平均、ロバストETS、負の二項、回帰付き時系列 &
大口注文が反復可能な需要か、一回限りの案件かを区別する。
外れ値を無条件に削除せず、案件情報があれば説明変数にする。\\
Intermittent &
季節ナイーブ、SBA、TSB、ハードルモデル &
需要発生確率と発生時数量を分ける。標準的な連続分布モデルだけで
ゼロを埋めない。\\
Lumpy &
TSB、ハードル・負の二項、品目間の階層モデル &
統計予測だけでは不確実性が大きい。大口案件、用途消滅、終売を
営業情報で確認し、高額品は人による承認対象とする。\\
\bottomrule
\end{tabular}
\end{center}

各クラスについて一つの「標準モデル」を先に固定するのではなく、
表の候補を同じローリング・オリジンで比較する。クラス内の全品目をまとめた
在庫費用、欠品率、MASE等から最良の方法をチャンピオン、次点を
チャレンジャーとする。個別品目で別モデルを採用するのは、十分な評価起点が
あり、改善が複数期間で再現した場合に限る。これにより、短い履歴への
過剰適合と品目ごとのモデル乱立を防ぐ。

データが少ない場合は、平滑化係数や過分散パラメータを品目ごとに推定せず、
同じクラスまたは同じ製品群で共有する方法も有効である。ただしクラス内でも
需要水準は異なるため、数量そのものを単純に平均せず、尺度調整または
階層モデルを用いる。

\subsection{事業情報と重要度を重ねる}

ADI--$\mathrm{CV}^2$ は販売履歴の形だけを表し、なぜその形になったかは
示さない。このため、導入・成長・成熟・終売、季節品、販促・大口案件、
用途消滅リスク、型番切替、欠品・供給制約のフラグを別に持つ。
これらは需要クラスを上書きするのではなく、モデルと人の判断を補足する。
\begin{itemize}
\item 新製品は履歴不足なのでADI分類をせず、類似品または製品群の情報と
販売計画を使う。実績が蓄積してから通常分類へ移す。
\item 旧品から新品への切替は二品目を独立に扱わず、合計需要と配分率を
予測する。旧品には終売日と残存サービス需要を設定する。
\item 欠品月の販売数は真の需要を下回るため、通常のゼロ・低需要として
学習させない。受注残や失注記録で補正するか、供給制約月として除外する。
\item 販促・大口案件は再発するなら説明変数、再発しないなら一時イベント
として記録する。単なる外れ値として処理すると将来の案件需要を見落とす。
\item ゼロが連続してIntermittentまたはLumpyへ移った品目は、
需要減少だけでなく用途消滅や後継品移行の確認対象とする。
\end{itemize}

さらに需要金額または粗利によるABC区分、顧客への重要度、欠品損失、
保管費用、陳腐化リスクを重ねる。例えばAランクかつLumpyの品目は
営業確認と個別承認を行い、CランクかつSmoothの品目は標準モデルで
自動処理する。サービス水準 $\alpha$ は需要クラスだけで決めず、
欠品費用と在庫費用、顧客重要度から決める。同じLumpyでも、欠品損失が
大きい品目と陳腐化費用が大きい品目では目標分位点が異なる。

\subsection{評価、製造指示、監視への活用}

モデル評価は全品目の平均だけでなく、クラス別に次を報告する。
\begin{itemize}
\item 品目数、需要金額、ゼロ需要率、ADIと $\mathrm{CV}^2$ の分布
\item 点予測のMASEとバイアス、予測区間の被覆率
\item $I+L$ 期間累計需要の分位点精度
\item 欠品率、充足率、平均在庫、過剰在庫、陳腐化費用
\end{itemize}
例えば全体精度が改善してもLumpy品の欠品が増えるなら、そのクラスだけ
モデルまたはサービス水準を見直せる。逆にSmooth品で複雑なモデルが
単純法を上回らなければ、単純法へ戻して運用を軽くできる。

製造指示では、分類そのものから製造本数を決めるのではない。
分類に応じて選んだモデルから $I+L$ 期間の累計需要分布を作り、
第\ref{sec:production-quantity}節の方法で目標分位点と
在庫ポジションとの差を製造数量とする。
分類は「どの予測分布を使うか」を決め、サービス水準と在庫状況は
「その分布のどこまで製造するか」を決める、という役割分担になる。

\subsection{再分類の運用ルール}

毎月指標を更新しても、境界を一度越えただけでモデルを即時変更すると
製造指示が不安定になる。次のようなルールを事前に定める。
\begin{enumerate}
\item ADI、$\mathrm{CV}^2$、ゼロ連続月数を毎月計算し、クラス判定は
四半期ごと等の一定周期で行う。
\item 新クラスの条件が2回連続の判定で成立した場合に移行する等、
ヒステリシスを設ける。境界付近には保留帯を置いてもよい。
\item クラスが変わっても直ちに本番モデルを替えず、新旧モデルを
並行評価し、在庫成果の改善を確認して切り替える。
\item SmoothからIntermittentへの移行、ゼロの長期継続、需要量の急増等は
統計的な再分類に加え、営業・商品担当へアラートを出す。
\item 各月について、分類値、選択モデル、上書き理由、製造指示量を保存し、
後から意思決定を再現できるようにする。
\end{enumerate}

導入時は「需要クラス $\times$ ABC区分」の表を作り、各セルについて
候補モデル、サービス水準の決定責任者、自動化可否、見直し周期を定める。
この運用表まで作ることで、需要分類を予測分析だけでなく実際の製造指示と
例外管理へ結び付けられる。

\section{評価方法}
\subsection{ローリング・オリジン設計}
過去への当てはまりでなく、実運用を再現するローリング・オリジン評価を
行う\cite{tashman2000}。各評価起点 $T_j$ で、その時点までに利用可能な
データだけを用いて、次を実施する。
\begin{enumerate}
\item 需要クラスと事業フラグを判定する。
\item 各候補モデルを再推定し、$T_j+1$ から $T_j+H$ の需要経路を生成する。
\item $D_{T_j,H}$ の分布と製造指示量を記録する。
\item 実現した需要と在庫推移を用いて予測・在庫成果を測る。
\end{enumerate}
起点を原則1か月ずつ進め、季節の全局面と複数の需要状態を含める。
最初の期間をモデル・係数・分類境界の選択、最後の連続した期間を最終評価に
用いる。最終評価を見て再調整すると成績を過大評価するため、調整後は
評価期間を固定する。比較対象は現行ガンマ法、単純法、パラメータ外挿法、
各クラスの候補モデルとする。

\subsection{予測精度の指標}
点予測誤差を $e_{T,h}=y_{T+h}-\hat y_{T+h|T}$ とする。
過小予測を正とするこの符号で平均誤差を報告すれば、正のバイアスは
欠品方向を示す。ゼロ需要では百分率誤差が未定義または不安定になるため、
MAPEを主指標にしない。
\begin{center}\begin{tabular}{>{\raggedright\arraybackslash}p{25mm}p{110mm}}\toprule
観点 & 指標\\\midrule
点予測 & MAE、バイアス、品目間比較が可能なMASE\cite{hyndman2006}\\
分布予測 & CRPSまたは対数スコア、50・80・95\%予測区間の被覆率と幅\cite{gneiting2007}\\
累計需要 & $H=I+L$ 期間累計のMAE、バイアス、分位点損失\\
在庫運用 & 欠品率、充足率、平均・過剰在庫、陳腐化、在庫費用と欠品費用\\
安定性 & 需要クラス別、需要額ABC別、ライフサイクル別成績、指示変更量\\\bottomrule
\end{tabular}\end{center}

MASEの尺度には、学習期間内の1期ナイーブ誤差または季節ナイーブ誤差を
用い、どちらを使ったかを固定する。分布予測では、名目95\,\%区間に
実績が概ね95\,\%入るかという校正と、区間が必要以上に広くないかという
鋭さを同時に確認する。製造指示は分位点に依存するため、全分布のCRPSに
加えて、実際に使用するサービス水準 $\alpha$ の分位点損失を重視する。

品目数だけの単純平均では低額品が結果を支配し、需要量による加重平均では
低需要品の問題が隠れる。したがって、品目単純平均、需要金額加重平均、
クラス別結果を併記する。評価起点または品目を単位とするブートストラップで
モデル間の差の区間を示し、僅かな順位差を確定的な改善と解釈しない。

\subsection{在庫シミュレーションと採用基準}
予測精度が高くても、製造ロット、リードタイム、既発行指示を無視すれば
在庫成果は改善しない。各評価起点で第\ref{sec:production-quantity}節
の規則を適用し、実際の $I,L$、初期在庫、製造残、ロット $B$、欠品時の
バックオーダーまたは失注規則を用いて在庫推移を再現する。
需要予測モデルとサービス水準 $\alpha$ は別々の意思決定変数なので、
まずモデルの分布校正を確認し、その後に費用または目標充足率に基づいて
$\alpha$ を選ぶ。

採用基準は評価前に定める。例えば「現行法と同等以上の充足率を保ちながら
平均在庫を削減し、かつ重要品クラスの欠品を悪化させない」とする。
平均成績だけでなく、最悪期の欠品、過剰在庫、指示量の月次変化、
担当者による上書き回数も制約に含める。

\section{適用上の限界}

本稿の方法を実装しても、次の不確実性は残る。
\begin{itemize}
\item 販売実績は供給制約下の潜在需要を観測しない。補正データがなければ、
どのモデルも需要を過小推定し得る。
\item 月次集計では月内の発生順序を失う。$I$ または $L$ が1か月より短い
場合は、週次・日次データで再構成する必要がある。
\item 新製品、突然の用途消滅、大口案件等、履歴にない構造変化は
時系列だけでは予測できず、営業・商品情報を必要とする。
\item 多数の候補から同じ評価期間で最良モデルを選ぶと過剰適合が起こる。
候補数を制限し、最終評価期間を分離する必要がある。
\item シミュレーション分布には、将来誤差だけでなく係数・状態・説明変数・
リードタイムの不確実性を可能な範囲で含める。点推定したパラメータを
固定すると予測区間が狭くなりやすい。
\item ガンマ近似は既存システムとの互換策であり、ゼロが多い分布や
多峰分布の分位点を正確に再現する保証はない。
\end{itemize}
したがって、予測結果は自動的な正解ではなく、再現可能な標準判断と
例外アラートを提供するものと位置付ける。

\section{実証および導入手順}
\begin{enumerate}
\item \textbf{データ監査：} 販売、欠品、受注残、価格、販促、切替関係、
在庫、製造残、$I,L$ を整理し、各項目の利用可能時点を記録する。
\item \textbf{探索：} 需要系列、移動統計、ADI、$\mathrm{CV}^2$、
ライフサイクルを可視化する。最低24--36か月、季節性には複数周期が望ましい。
\item \textbf{小規模実証：} 全クラスとABC区分を含む代表品目で、現行法、
単純法、外挿法、ETS、カウント、SBA/TSB、ハードル法を比較する。
\item \textbf{標本外評価：} 固定したローリング評価設計で分布校正と
$H$ 期間累計需要を比較し、クラスごとの候補を絞る。
\item \textbf{在庫評価：} 実際の指示周期、リードタイム、ロット制約を
含むシミュレーションを行い、採用基準を満たすか判定する。
\item \textbf{並行稼働：} 3--6か月は現行方式と新方式を並行運用し、
指示差、上書き理由、実績を保存する。
\item \textbf{本番監視：} バイアス、区間被覆率、充足率、在庫、
クラス移動、データ欠損を毎月監視し、所定周期で再推定する。
\end{enumerate}

\section{結論}
12か月窓の情報が遅れて需要変化への追随性を損なうという問題設定は妥当で
ある。ただし、窓推定値を約5.5か月外挿すれば現在の真の分布が得られるとは
限らない。窓の重複、少数標本、ゼロ需要による不安定性が残るためである。

実務上は、パラメータ外挿を現行法から移行しやすいベースラインとし、
月次需要を直接扱う状態空間・カウント・ハードルモデルを本命候補とする。
断続需要にはSBA、TSBも含める。ADI--$\mathrm{CV}^2$分類はモデルを
自動決定する規則ではなく、候補の絞り込み、評価の層別化、例外管理に用いる。
最終的な製造数量は、採用モデルから得た $I+L$ 期間累計需要の分位点と
在庫ポジションとの差として決める。

提案手法の採否は、同一条件のローリング・オリジン評価と在庫シミュレーション
で判断する。予測誤差だけでなく、分布の校正、充足率、平均在庫、陳腐化、
指示の安定性を評価し、現行方式に対する業務上有意な改善を確認して導入する。

\section{参考文献（概要付）}
\begin{thebibliography}{99}
\bibitem{croston1972} J. D. Croston, Forecasting and Stock Control for Intermittent Demands,
\textit{Operational Research Quarterly}, 23(3), 289--303, 1972.
\url{https://doi.org/10.1057/jors.1972.50}

正の需要量と需要間隔を別々に平滑化する方法を提案した基礎論文。

\bibitem{syntetos2005accuracy} A. A. Syntetos and J. E. Boylan, The Accuracy of Intermittent Demand Estimates,
\textit{International Journal of Forecasting}, 21(2), 303--314, 2005.
\url{https://doi.org/10.1016/j.ijforecast.2004.10.001}

Croston法の偏りを分析し、補正法SBAを検討した。

\bibitem{syntetos2005category} A. A. Syntetos, J. E. Boylan and J. D. Croston,
On the Categorization of Demand Patterns,
\textit{Journal of the Operational Research Society}, 56(5), 495--503, 2005.
\url{https://doi.org/10.1057/palgrave.jors.2601841}

ADIと正の需要量の変動による分類を理論的誤差比較から構成した。

\bibitem{teunter2011} R. H. Teunter, A. A. Syntetos and M. Z. Babai,
Intermittent Demand: Linking Forecasting to Inventory Obsolescence,
\textit{European Journal of Operational Research}, 214(3), 606--615, 2011.
\url{https://doi.org/10.1016/j.ejor.2011.05.018}

発生確率を毎期更新するTSB法を提案し、終売・陳腐化への対応を示した。

\bibitem{tashman2000} L. J. Tashman, Out-of-sample Tests of Forecasting Accuracy,
\textit{International Journal of Forecasting}, 16(4), 437--450, 2000.
\url{https://doi.org/10.1016/S0169-2070(00)00065-0}

標本外評価について、移動起点、窓、再推定等の設計を整理した。

\bibitem{hyndman2006} R. J. Hyndman and A. B. Koehler, Another Look at Measures of Forecast Accuracy,
\textit{International Journal of Forecasting}, 22(4), 679--688, 2006.
\url{https://doi.org/10.1016/j.ijforecast.2006.03.001}

誤差指標の問題を整理し、尺度に依存しないMASEを提案した。

\bibitem{gneiting2007} T. Gneiting and A. E. Raftery, Strictly Proper Scoring Rules, Prediction, and Estimation,
\textit{Journal of the American Statistical Association}, 102(477), 359--378, 2007.
\url{https://doi.org/10.1198/016214506000001437}

CRPS等の適正スコアを体系化し、確率分布全体を評価する根拠を与えた。

\bibitem{hyndman2002ets}
R. J. Hyndman, A. B. Koehler, R. D. Snyder and S. Grose,
A State Space Framework for Automatic Forecasting Using Exponential Smoothing Methods,
\textit{International Journal of Forecasting}, 18(3), 439--454, 2002.
\url{https://doi.org/10.1016/S0169-2070(01)00110-8}

指数平滑法を状態空間モデルとして統一し、尤度による選択、予測区間、
将来経路のシミュレーションを可能にした。

\bibitem{davis2003count}
R. A. Davis, W. T. M. Dunsmuir and S. B. Streett,
Observation-driven Models for Poisson Counts,
\textit{Biometrika}, 90(4), 777--790, 2003.
\url{https://doi.org/10.1093/biomet/90.4.777}

過去の観測と説明変数に依存するPoissonカウント時系列モデルを定式化し、
推定量の性質と適用例を示した。

\bibitem{mullahy1986}
J. Mullahy, Specification and Testing of Some Modified Count Data Models,
\textit{Journal of Econometrics}, 33(3), 341--365, 1986.
\url{https://doi.org/10.1016/0304-4076(86)90002-3}

ゼロの発生と正のカウントを分ける修正カウントモデルを提示した、
ハードル型モデルの基礎文献である。
\end{thebibliography}
\end{document}
